\begin{abstract}
Conditional invariance penalties are commonly used to improve EEG domain generalization, but global
penalties can remove task information, collapse representations, or optimize a leakage proxy without
improving target performance. We study a selective alternative: localize conditional domain information in
task-light EEG subspaces using score-Fisher geometry, and intervene only when deletion passes an explicit
source-risk certificate.

Synthetic controls show why certification is necessary. First-moment leakage proxies miss covariance-mediated
domain information, geometry-only task protection can still delete conditionally task-useful information, and
weak learned gates can unsafe-accept Bayes-unsafe deletions. A cross-fitted log-ratio task-risk gate reduces
this failure mode, but at moderate sample sizes the certified action is often refusal rather than deletion.

On BCI-IV-2a, the framework reveals a measurement-to-control gap. In TSMNet, subject identity is localizable
but redundantly encoded across a high-dimensional latent space: low-rank deletion preserves the task but only
dents leakage, while global LPC removes leakage only by collapsing the representation to the origin. In
EEGNet, the same diagnostic deletion removes much more subject leakage with negligible task cost, yet target
accuracy does not improve. These results suggest that conditional domain leakage is measurable, and sometimes
removable, but neither measurement nor removal is sufficient evidence of domain-generalization benefit. Safe
selective invariance should therefore be treated as a certified intervention with refusal, not as an
always-on regularizer.
\end{abstract}
